\scanpage{43}\mspage{42}
\uline{caractérisée} par les conditions suivantes :
% slanted
\margin{C'est faux, ça implique (i) (ii) (iii) (iv), mais ces conditions ne suffisent nullement, voir plus bas ! v.~48}
\margin{i.e $\overline{\tau'}\in\underline{\underline{\overline{\mathscr{T}}}}$ $\left\{\vphantom{\rule{0pt}{4ex}}\right.$}

(i) $\overline{\tau'}$ est plus fine que $\widetilde{\tau}$, et induit sur $J'$ $\widetilde{\tau}\mid J'$.

(ii) $\overline{\tau'}\mid\overset{\deleted{\unsure{\sim|\mathord{-}}}}{U'}=I'\setminus J'$ coïncide avec $\widetilde{\tau}\mid\overset{\deleted{\sim}}{U'}$

(iii) $\varphi\mid U'$ induit \deleted{une bijection} \inserted{\unsure{iso. préord.}} $U'\longrightarrow U=S\setminus T$
% slanted
\margin{\deleted{Autant dire}}

(\unsure{iv|iiv}) La topologie de $I$ est quotient final pour les applications $I'\mathbin{\text{---}}I$, $J\mathbin{\text{---}}I$ ($I'$ munie \deleted{celle de $I$, \unsure{conjointement}} \inserted{de $\overline{\tau'}$, $J$ \deleted{\unsure{de}} muni de $\overline{\tau}\mid J$}) (Ce qui exclut notamment la solution triviale pour (i) (ii), obtenue en prenant la topologie somme sur \deleted{$J'$} $I'=J'\underset{\text{disj.}}{\cup}\overset{\deleted{\sim}}{U'}$, sauf si $I$ lui$=$ est une somme top. $J\sqcup\underset{\substack{\Vert\\I\setminus J}}{U}$).
% slanted
\margin{C'est cette condition que j'avais d'abord oubliée, à la page 31 !}
% slanted
\margin{\deleted{Autant} dire que $\tau'$ sur $S$ est caractérisée par les propriétés analogues (i) (ii) (iii) (iv), la dernière \unsure{signifiant|signifie} que la top $\text{\textcircled{$\tau$}}$ sur $S$ est finale pour $\unsure{\tau''|\tau'}$ sur $S$ et $\tau$ sur $T$ (id $S\unsure{\Rightarrow|\rightrightarrows|\to}S$ $T\unsure{\hookrightarrow|\to}S$)}
% slanted
\margin{en fait, c'est indépendant aussi de l'hyp. \unsure{ultra-topologique}}

Démonstration

\deleted{a) et b).} \inserted{(a b c)} Il est clair que les ens. épinglés dans c) sont ceux qui satisfont $1^\circ)$, $2^\circ)$ $3^\circ)$ dans b), et que la représentation dans c) est unique. Il reste à voir, pour prouver (a,b,c), que les ens. forment une topologie, \deleted{et} qu'elle \deleted{\unsure{a les propriétés|a les propri.}} \inserted{est} dans $\underline{\mathscr{T}}$, et est moins fine des top $\in\underline{\mathscr{T}}$.

1) Ce sont les fermés d'une topologie.

Stables par intersections quelconques : clair

Stables par réunions finies : idem.

C'est $=$ stable par réunions quelconques s'il en est ainsi pour $(T,\tau'_0)$, $(\deleted{\unsure{I|D}}\unsure{S|I},\tau)$, d'où la première assertion de \unsure{d)|a)}.
